
\chapter{Proposed Methodology}
\label{chap:methodology}
\section{Methodology Overview}

Stage 1 showed that adversarial robustness does not transfer uniformly across
attack families. A model trained against one threat model can remain vulnerable
to attacks with different perturbation geometries, optimization procedures, or
loss objectives. This motivates the Stage 2 methodology: instead of treating all
adversarial examples as samples from one homogeneous distribution, this chapter
models attack-induced examples as domains and asks how training can emphasize
weak domains or latent hard groups without losing the stability of the uniform
multi-attack baseline.

The chapter develops this progression step by step. Section~\ref{sec:ch4_problem_formulation}
first formulates multi-attack adversarial training under the attacks-as-domains
perspective. Section~\ref{sec:ch4_multi_attack_erm} then defines Multi-AT, the
uniform multi-attack baseline used alongside the single-attack PGD-AT and DDN-AT
baselines. Section~\ref{sec:ch4_attackdro} introduces AttackDRO, which applies
group distributionally robust optimization (Group DRO) over source attack
identities, but this still assumes that each attack forms one homogeneous group.

The remaining sections address that limitation. Section~\ref{sec:ch4_attackdropp}
introduces AttackDRO++, which replaces attack identities with discovered
clusters. Section~\ref{sec:ch4_gradfp} augments those clusters with gradient
fingerprints so that grouping can reflect optimization-level difficulty, not
only representation similarity. Finally, Section~\ref{sec:ch4_anchor_objective}
adds a uniform anchor that keeps the final AttackDRO++ objective tied to the
Multi-AT training signal while still allowing cluster-level correction.

% --------------------------------------------------------------------------
\section{Multi-Attack Training as a Domain Problem}
\label{sec:ch4_problem_formulation}

\subsection{Setup and Notation}

Let
\begin{equation}
    \mathcal{D}_{\mathrm{train}}
    =
    \{(\mathbf{x}_i, y_i)\}_{i=1}^{n}
\end{equation}
denote the supervised training dataset, where \(\mathbf{x}_i \in \mathcal{X}\) is an
input image and \(y_i \in \{1,\ldots,C\}\) is its class label. Let \(f_\theta\) be a
classifier parameterized by \(\theta\), and let
\(\ell(f_\theta(\mathbf{x}),y)\) denote the cross-entropy loss.

We denote the logits produced by the model as
\[
    \mathbf{z} = f_\theta(\mathbf{x}) \in \mathbb{R}^{C},
\]
and the penultimate-layer representation as
\[
    \mathbf{h} = \phi_\theta(\mathbf{x}) \in \mathbb{R}^{d_h}.
\]
The representation \(\phi_\theta(\cdot)\) is later used for cluster discovery.

We define a set of source attack domains
\begin{equation}
    \mathcal{A}_{\mathrm{src}}
    =
    \{A_1,A_2,\ldots,A_M\},
\end{equation}
where each attack operator \(A_m\) is characterized by a threat model, perturbation
budget, optimization procedure, number of inner steps, and loss objective. Given a
clean sample \((\mathbf{x}_i,y_i)\), the adversarial example produced by source attack
\(A_m\) is
\begin{equation}
    \mathbf{x}_i^{(m)}
    =
    A_m(f_\theta,\mathbf{x}_i,y_i).
\end{equation}

In the main experiments, we use two source attacks: a PGD-\(\ell_\infty\) attack and
a DDN-\(\ell_2\) attack. This choice exposes the model to two common adversarial norm
families during training. Additional attacks are reserved as heldout evaluation attacks
and are never used for optimization.

\subsection{Multi-Attack Risk and Its Limitations}

The empirical adversarial risk under source attack \(A_m\) is
\begin{equation}
    R_m(\theta)
    =
    \frac{1}{n}
    \sum_{i=1}^{n}
    \ell\!\left(
        f_\theta\!\left(\mathbf{x}_i^{(m)}\right),
        y_i
    \right).
\end{equation}

A standard multi-attack objective minimizes the average risk across source attacks:
\begin{equation}
    \min_\theta
    \frac{1}{M}
    \sum_{m=1}^{M}
    R_m(\theta).
    \label{eq:avg_risk}
\end{equation}

This objective improves attack diversity compared with single-attack adversarial
training. However, it treats all source attacks and all adversarial examples with equal
importance. This creates two limitations.

First, if one source attack consistently induces higher loss than another, equal
weighting may underemphasize the harder attack. Second, adversarial examples within
the same attack are not homogeneous. Two samples generated by the same algorithm can
differ substantially in classification margin, representation, and optimization
direction. These within-attack differences are not captured by
Equation~\eqref{eq:avg_risk}.

The central goal of this chapter is therefore to construct adaptive objectives that can
identify and upweight hard groups while preserving the stability of the uniform
multi-attack training signal.

% ─────────────────────────────────────────────────────────────────────────────
\section{Uniform Multi-Attack ERM Baseline}
\label{sec:ch4_multi_attack_erm}

Uniform Multi-Attack empirical risk minimization (ERM) serves as the main baseline
for this work. It trains on multiple source attacks within the same minibatch but
assigns equal weight to all generated adversarial examples. This makes it a strong
and simple reference point for evaluating whether adaptive reweighting provides
additional benefit.

Given a minibatch
\[
    \mathcal{B} = \{(\mathbf{x}_i,y_i)\}_{i=1}^{N},
\]
we partition it into \(M\) disjoint subsets
\[
    \mathcal{B}_1,\ldots,\mathcal{B}_M,
    \qquad
    \mathcal{B}_m \cap \mathcal{B}_{m'} = \emptyset
    \;\; \text{for } m\ne m',
    \qquad
    \bigcup_{m=1}^{M}\mathcal{B}_m = \mathcal{B}.
\]
Each subset \(\mathcal{B}_m\) is assigned to one source attack \(A_m\). For each
\((\mathbf{x}_i,y_i)\in\mathcal{B}_m\), we generate
\[
    \mathbf{x}_i^{\mathrm{adv}}
    =
    A_m(f_\theta,\mathbf{x}_i,y_i).
\]

The adversarial examples from all source attacks are then concatenated into one mixed
adversarial batch. The uniform multi-attack loss is
\begin{equation}
    \mathcal{L}_{\mathrm{ERM}}
    =
    \frac{1}{N}
    \sum_{i=1}^{N}
    \ell\!\left(
        f_\theta(\mathbf{x}_i^{\mathrm{adv}}),
        y_i
    \right).
    \label{eq:erm_loss}
\end{equation}

In our main setting, the source attacks are PGD-\(\ell_\infty\) and DDN-\(\ell_2\),
so the model is exposed to both \(\ell_\infty\) and \(\ell_2\) perturbation
geometries during training. However, the objective in
Equation~\eqref{eq:erm_loss} treats all adversarial examples equally. It does not
explicitly identify which attack, class, or latent subgroup is harder. This limitation
motivates the adaptive reweighting methods introduced in the following sections.

% ─────────────────────────────────────────────────────────────────────────────
\section{AttackDRO: Group DRO Over Attack Identities}
\label{sec:attackdro}
% ─────────────────────────────────────────────────────────────────────────────

AttackDRO extends Uniform Multi-Attack ERM by applying Group Distributionally Robust
Optimization (GroupDRO) over attack identities. Each source attack \(A_m\) is treated
as one group. The method maintains an attack-group weight vector
\begin{equation}
    \mathbf{q} = (q_1,\ldots,q_M) \in \Delta^M,
\end{equation}
where \(\Delta^M\) is the \(M\)-dimensional probability simplex.

For each attack group \(m\), the group loss is
\begin{equation}
    \mathcal{L}_m
    =
    \frac{1}{|\mathcal{B}_m|}
    \sum_{i \in \mathcal{B}_m}
    \ell\!\left(
        f_\theta(\mathbf{x}_i^{\mathrm{adv}}),
        y_i
    \right).
\end{equation}
The AttackDRO objective is then
\begin{equation}
    \mathcal{L}_{\mathrm{AttackDRO}}
    =
    \sum_{m=1}^{M} q_m \mathcal{L}_m.
\end{equation}

The group weights are updated using an exponentiated-gradient (EG) rule. Let \(s_m\)
be the difficulty score of attack group \(m\). The update is
\begin{equation}
    q_m \leftarrow q_m \exp(\eta_q s_m),
    \qquad m=1,\ldots,M,
    \label{eq:eg}
\end{equation}
followed by normalization and floor projection:
\begin{equation}
    q_m \ge q_{\min},
    \qquad
    \sum_{m=1}^{M} q_m = 1.
    \label{eq:attackdro_floor}
\end{equation}
The floor constraint prevents any attack group from being completely ignored.

In our formulation, the difficulty score is based on the classification margin. For a
sample \(i\) with logits \(\mathbf{z}_i\), define the margin as
\begin{equation}
    r_i
    =
    z_{i,y_i}
    -
    \max_{c \ne y_i} z_{i,c}.
\end{equation}
A large positive margin indicates a confident correct prediction, while a small or
negative margin indicates a hard or misclassified sample. We convert this margin into
a per-sample difficulty score using
\begin{equation}
    s_i
    =
    \mathrm{softplus}(-r_i/\tau)
    =
    \log(1+\exp(-r_i/\tau)),
\end{equation}
where \(\tau>0\) is a temperature parameter. The group-level difficulty score is the
average difficulty within group \(m\):
\begin{equation}
    s_m
    =
    \frac{1}{|\mathcal{B}_m|}
    \sum_{i\in\mathcal{B}_m} s_i.
\end{equation}

AttackDRO can adaptively emphasize whichever source attack is currently harder.
However, it relies on a coarse grouping assumption: all examples generated by the same
attack are treated as one homogeneous group. In practice, two PGD-\(\ell_\infty\)
examples may have very different margins, representations, and gradient directions.
Conversely, examples generated by different attacks may share similar failure modes.
Thus, the EG update can only adapt at the attack level and may miss finer-grained
difficulty structure within and across attacks. The cluster-based extension in the next
section addresses this limitation.

% --------------------------------------------------------------------------
\section{Cluster-Level DRO: Discovering Latent Groups}
\label{sec:ch4_attackdropp}

This section introduces the cluster-level DRO component used by AttackDRO++. Instead
of assigning one group to each source attack identity, the component discovers latent
groups from adversarial examples. The central hypothesis is that attack identity is
too coarse as a grouping criterion. Two adversarial examples generated by different
attacks may share the same failure mode, while two examples generated by the same
attack may have very different margins, representations, and optimization directions.
Cluster-based grouping allows the DRO objective to adapt to hard subgroups that may
cut across attack boundaries. The complete AttackDRO++ method is defined after the
augmented clustering features and uniform anchor are introduced in Sections~\ref{sec:ch4_gradfp}
and~\ref{sec:ch4_anchor_objective}.

For each adversarial example \(\mathbf{x}_i^{\mathrm{adv}}\), the model extracts a
representation vector
\begin{equation}
    \mathbf{h}_i = \phi_\theta(\mathbf{x}_i^{\mathrm{adv}}),
\end{equation}
and a clustering model assigns the example to a latent group
\begin{equation}
    c_i \in \{1,\ldots,K\}.
\end{equation}
Let \(\mathcal{B}_k\) denote the subset of the current minibatch whose adversarial
examples are assigned to cluster \(k\):
\begin{equation}
    \mathcal{B}_k
    =
    \{i \in \{1,\ldots,N\}: c_i = k\}.
\end{equation}
The cluster loss is
\begin{equation}
    \mathcal{L}_k
    =
    \frac{1}{|\mathcal{B}_k|}
    \sum_{i\in\mathcal{B}_k}
    \ell\!\left(
        f_\theta(\mathbf{x}_i^{\mathrm{adv}}),
        y_i
    \right).
    \label{eq:cluster_loss}
\end{equation}

The cluster-DRO component maintains a cluster-level weight vector
\begin{equation}
    \mathbf{q} = (q_1,\ldots,q_K) \in \Delta^K.
\end{equation}
Because not every cluster appears in every minibatch, the cluster weights are
renormalized over the clusters present in the current minibatch. Let
\begin{equation}
    \mathcal{C}_{\mathcal{B}}
    =
    \{k \in \{1,\ldots,K\}: |\mathcal{B}_k|>0\}
\end{equation}
denote the set of present clusters. For each \(k \in \mathcal{C}_{\mathcal{B}}\), we
define
\begin{equation}
    \tilde{q}_k
    =
    \frac{q_k}
    {\sum_{j\in\mathcal{C}_{\mathcal{B}}} q_j}.
\end{equation}
The cluster-DRO loss is then
\begin{equation}
    \mathcal{L}_{\mathrm{ClusterDRO}}
    =
    \sum_{k\in\mathcal{C}_{\mathcal{B}}}
    \tilde{q}_k \mathcal{L}_k.
    \label{eq:cluster_dro_loss}
\end{equation}

This objective changes the unit of robustness balancing. Rather than assigning one
weight to each attack identity, the model assigns weight to discovered subgroups
of adversarial examples, so hard regions that cut across attack families can receive
more training emphasis. The renormalization over present clusters keeps the minibatch
objective well-defined even when some latent groups are absent from a particular batch.

The cluster weights are updated using the same exponentiated-gradient principle as in
AttackDRO, but now at the cluster level rather than the attack level. Given a
cluster-level difficulty score
\begin{equation}
    s_k
    =
    \frac{1}{|\mathcal{B}_k|}
    \sum_{i\in\mathcal{B}_k} s_i,
    \qquad k \in \mathcal{C}_{\mathcal{B}},
\end{equation}
the update increases the weight of clusters with higher average difficulty. This allows
the cluster-DRO component to emphasize hard latent groups even when they do not align exactly with
the original attack identities.

Since model representations evolve during training, cluster assignments must also be
updated. The clustering procedure therefore maintains a feature bank \(\mathcal{F}\), which stores
detached augmented feature vectors from recent minibatches. When the refresh condition
is met, K-means is refitted on \(\mathcal{F}\), and subsequent minibatches are assigned
using the updated clusterer. Before the first K-means refresh, cluster assignments are
bootstrapped randomly.

The refresh condition is treated as a policy rather than a fixed part of the method. It
can be defined in terms of epochs or training steps. In the main experiments, we use an
epoch-based refresh schedule with \(R=2\) epochs. After each refresh, the feature bank
is reset or updated according to the memory policy, and the cluster weights are refreshed
according to the chosen \(q\)-refresh policy.

% --------------------------------------------------------------------------
\section{Augmented Clustering with Gradient Fingerprints}
\label{sec:ch4_gradfp}

Representation-space clustering alone may not fully capture adversarial difficulty.
Two adversarial examples can be close in feature space while inducing different
optimization behavior. To capture this additional signal, AttackDRO++ constructs
an augmented clustering feature that combines representation similarity, class
information, and gradient-level training dynamics.

For an adversarial example \(\tilde{\mathbf{x}}_i=\mathbf{x}_i^{\mathrm{adv}}\), let
\[
z_i = f_\theta(\tilde{\mathbf{x}}_i),
\qquad
h_i = \phi_\theta(\tilde{\mathbf{x}}_i).
\]
The logits \(z_i=f_\theta(\tilde{\mathbf{x}}_i)\) are used to compute prediction
probabilities and gradient-fingerprint terms, while
\(h_i=\phi_\theta(\tilde{\mathbf{x}}_i)\) provides the representation component of
the clustering feature. Let
\[
p_i = \mathrm{softmax}(z_i)
\]
be the predicted class distribution.

The first component is the normalized representation feature,
\[
\bar{h}_i = \operatorname{norm}(h_i),
\]
which captures semantic similarity in the model feature space.

The second component is a class-aware scalar label feature. When label information
is enabled, the implementation appends the normalized scalar class-index feature
\(\tilde{y}_i \in [0,1]\) with weight \(\lambda_{\mathrm{lbl}}\):
\[
\lambda_{\mathrm{lbl}}\tilde{y}_i,
\]
which provides weak label structure during clustering without using a one-hot label
vector in the clustering feature.

The third component is the gradient fingerprint. For a linear classifier head
\[
f_\theta(\cdot) = W\phi_\theta(\cdot) + b,
\qquad
W \in \mathbb{R}^{C \times d_h},
\]
the per-example gradient of the cross-entropy loss with respect to the classifier
weight matrix is
\[
\nabla_W \ell(f_\theta(\tilde{\mathbf{x}}_i),y_i)
=
\bigl(p_i-\mathrm{onehot}(y_i)\bigr)\otimes h_i
\in \mathbb{R}^{C\times d_h}.
\]
We use this quantity as a gradient proxy and denote it by
\[
G_i =
\bigl(p_i-\mathrm{onehot}(y_i)\bigr)\otimes h_i.
\]
The one-hot vector appears here only in the cross-entropy gradient proxy; it is
not appended as the label component of \(\psi_i\).

Intuitively, \(G_i\) describes the training direction induced by sample \(i\).
Thus, two adversarial examples may be grouped together not only because they are
close in representation space, but also because they induce similar optimization
updates.

The raw gradient proxy has dimension \(C d_h\). For CIFAR-10 with a ResNet-18
backbone, this corresponds to \(10 \times 512 = 5120\) dimensions. To make the
fingerprint more compact and efficient for clustering, we use a fixed random
projection matrix
\[
P \in \mathbb{R}^{C d_h \times d_{\mathrm{proj}}},
\]
and define the projected gradient fingerprint as
\[
g_i = \mathrm{vec}(G_i)P
\in \mathbb{R}^{d_{\mathrm{proj}}}.
\]
In our default setting, \(d_{\mathrm{proj}}=128\). The projection is fixed throughout
training and introduces no additional learnable parameters.

The final augmented clustering feature is
\[
\psi_i =
\left[
\operatorname{norm}(h_i)
\;\Vert\;
\lambda_{\mathrm{lbl}}\tilde{y}_i
\;\Vert\;
\lambda_{\mathrm{grad}}\operatorname{norm}(g_i)
\right].
\]

The three components play complementary roles. Representation features describe where
an example lies in feature space, whereas gradient fingerprints describe how the
classifier head would need to change to classify that example correctly. Combining
representation, scalar label, and gradient information makes the clustering signal
both representation-aware and optimization-aware. This allows AttackDRO++ to discover
latent adversarial groups that reflect both structural similarity and optimization
difficulty.

% --------------------------------------------------------------------------
\section{Uniform-Anchored Training Objective}
\label{sec:ch4_anchor_objective}

A pure Cluster-DRO objective can become unstable when the cluster weights
\(\mathbf{q}\) concentrate too strongly on a small number of hard clusters. We refer
to this failure mode as \emph{\(q\)-collapse}. In the extreme case, the objective
approaches a worst-cluster loss, where most of the training signal is dominated by
one cluster. This behavior is undesirable in multi-attack adversarial training because
the hardest cluster may change after each cluster refresh, and over-emphasizing one
cluster can reduce robustness on the remaining attacks or groups.

To stabilize training, AttackDRO++ uses a uniform-anchored Cluster-DRO objective
that interpolates between the uniform multi-attack ERM loss and the adaptive
cluster-DRO loss:
\begin{equation}
    \mathcal{L}_{\mathrm{AttackDRO++}}
        =
        (1 - \lambda_{\mathrm{DRO}})\mathcal{L}_{\mathrm{ERM}}
        +
        \lambda_{\mathrm{DRO}}\mathcal{L}_{\mathrm{ClusterDRO}} .
    \label{eq:final_loss}
\end{equation}
Here, \(\lambda_{\mathrm{DRO}}\in[0,1]\) controls the strength of the Cluster-DRO
correction. When \(\lambda_{\mathrm{DRO}}=0\), the method reduces to uniform
multi-attack ERM. When \(\lambda_{\mathrm{DRO}}=1\), the method becomes pure
Cluster-DRO.

Equivalently, Equation~\eqref{eq:final_loss} can be written as
\begin{equation}
    \mathcal{L}_{\mathrm{AttackDRO++}}
        =
        \mathcal{L}_{\mathrm{ERM}}
        +
        \lambda_{\mathrm{DRO}}
        \left(
            \mathcal{L}_{\mathrm{ClusterDRO}}
            -
            \mathcal{L}_{\mathrm{ERM}}
        \right).
\end{equation}
This form shows that \(\lambda_{\mathrm{DRO}}\) scales the difference between the
cluster-reweighted loss and the uniform loss. If the cluster-DRO loss is close to the
uniform loss, the correction is small. If some clusters are substantially harder, the
correction becomes larger, but its influence remains controlled by
\(\lambda_{\mathrm{DRO}}\).

In the main experiments, we use \(\lambda_{\mathrm{DRO}}=0.35\), used as the default anchor strength in the main configuration. This value keeps the uniform multi-attack objective as
the dominant stabilizing term while still allowing the model to adapt toward harder
clusters. Intuitively, this suggests an inverted-U behavior: a small amount of
Cluster-DRO pressure can improve robustness by emphasizing hard groups, whereas
excessive pressure may overfit unstable clusters and harm the balance across attacks.

A second stabilization mechanism is the floor projection applied to the cluster
weights. After each exponentiated-gradient update, \(\mathbf{q}\) is normalized and
projected so that every cluster keeps at least \(q_{\min}\) mass:
\begin{equation}
    q_k \ge q_{\min}, \qquad \sum_{k=1}^{K} q_k = 1.
    \label{eq:floor}
\end{equation}

The anchor controls how much the adaptive cluster loss can influence the total objective,
while the floor projection controls the cluster weights themselves. Together, they
prevent the adaptive component from collapsing onto a small number of clusters and keep
all discovered groups represented during training. In our default setting,
\(q_{\min}=0.03\).

Finally, the \(q\)-update is delayed during the early training phase. For the first
\(T_w\) epochs, \(\mathbf{q}\) remains uniform while the model representation and
initial cluster assignments become more reliable. After this warmup period, the
cluster weights are updated using the exponentiated-gradient rule. In the main
configuration, \(T_w=3\) epochs.

% --------------------------------------------------------------------------
\section{Complete Training Framework}
\label{sec:ch4_complete_pipeline}

AttackDRO++ is the final method studied in this report. It combines latent
cluster discovery, augmented clustering features, gradient fingerprints, and a
uniform-anchored Cluster-DRO objective into one training pipeline. The goal is to
move beyond fixed attack identities and let the model identify harder adversarial
subgroups during training. At the same time, the uniform anchor prevents the
adaptive objective from drifting too far from the strong multi-attack baseline.
Algorithm~\ref{alg:ch4_attackdropp} gives the complete procedure, while
Table~\ref{tab:ch4_hparams} lists the default hyperparameters used in the main
experiments.

\paragraph{From Diagnostic Motivation to the Final Method}

The method follows directly from the attacks-as-domains diagnosis introduced
earlier. Single-attack adversarial training reveals the cross-attack robustness
gap: a model trained against one source attack can remain weak against other
threat models. Uniform multi-attack training addresses this by exposing the model
to multiple source attacks, but it treats all generated examples equally.
AttackDRO then adds attack-level reweighting, but this still assumes that each
attack family is internally homogeneous.

AttackDRO++ removes this assumption. Instead of using attack identity as the only
group label, it discovers latent groups from the adversarial examples themselves.
Gradient fingerprints add optimization-level information to the clustering
features, helping the grouping capture not only representation similarity but also
training difficulty. Finally, the uniform anchor keeps the objective close to the
strong multi-attack baseline while still allowing Cluster-DRO to emphasize harder
groups. The result is a method that moves from observing the cross-attack gap to
actively training against the latent structure behind it.

The same evaluation discipline is preserved throughout the framework. Source
attacks are fixed during training, held-out attacks are never used for
optimization, and robustness is reported both per attack and through aggregate
metrics.

The dominant cost remains adversarial example generation, which is shared by all
multi-attack methods in the comparison. Therefore, AttackDRO++ mainly changes the
training objective and grouping mechanism, while keeping the source attacks,
threat models, and evaluation protocol fixed. The cluster refresh schedule is an implementation choice rather than a fixed part
of the mathematical objective. In the main experiments, clusters are refreshed
every \(R=2\) epochs, and \(\mathbf{q}\) is reset to the uniform distribution
after each refresh to avoid carrying weights across changed cluster identities.


\begin{figure}[htbp]
\centering
\resizebox{0.95\textwidth}{!}{%
\begin{tikzpicture}[
  node distance = 0.68cm,
  every node/.append style = {align=center},
  inputblk/.style  = {draw=black!40,        fill=gray!11,
                      rounded corners=5pt,  minimum width=7.6cm,
                      minimum height=1.08cm, font=\small,
                      inner sep=8pt,        line width=0.85pt},
  attackblk/.style = {draw=red!50!black,    fill=red!9,
                      rounded corners=5pt,  minimum width=7.6cm,
                      minimum height=1.08cm, font=\small,
                      inner sep=8pt,        line width=0.85pt},
  fwdblk/.style    = {draw=blue!55!black,   fill=blue!8,
                      rounded corners=5pt,  minimum width=7.6cm,
                      minimum height=1.08cm, font=\small,
                      inner sep=8pt,        line width=0.85pt},
  featblk/.style   = {draw=violet!65!black, fill=violet!8,
                      rounded corners=5pt,  minimum width=7.6cm,
                      minimum height=1.08cm, font=\small,
                      inner sep=8pt,        line width=0.85pt},
  clustblk/.style  = {draw=orange!75!black, fill=orange!10,
                      rounded corners=5pt,  minimum width=7.6cm,
                      minimum height=1.08cm, font=\small,
                      inner sep=8pt,        line width=0.85pt},
  lossblk/.style   = {draw=green!60!black,  fill=green!9,
                      rounded corners=5pt,  minimum width=3.5cm,
                      minimum height=1.60cm, font=\small,
                      inner sep=8pt,        line width=0.85pt},
  combblk/.style   = {draw=teal!65!black,   fill=teal!9,
                      rounded corners=5pt,  minimum width=7.6cm,
                      minimum height=1.08cm, font=\small,
                      inner sep=8pt,        line width=0.85pt},
  updblk/.style    = {draw=teal!80!black,   fill=teal!13,
                      rounded corners=5pt,  minimum width=7.6cm,
                      minimum height=1.08cm, font=\small,
                      inner sep=8pt,        line width=0.85pt},
  sideblk/.style   = {draw=gray!55,         fill=gray!8,
                      rounded corners=5pt,  minimum width=3.1cm,
                      minimum height=1.05cm, font=\small,
                      inner sep=7pt,        line width=0.72pt,
                      densely dashed},
  arr/.style       = {-{Stealth[length=7pt, width=5.5pt]},
                      line width=0.95pt, color=black!55},
  darr/.style      = {-{Stealth[length=6pt, width=4.5pt]},
                      line width=0.72pt, color=gray!60, densely dashed},
  looparr/.style   = {-{Stealth[length=7pt, width=5.5pt]},
                      line width=0.80pt, color=black!28, rounded corners=4pt},
]

\node[inputblk] (batch)
  {\textbf{Minibatch}\qquad
   $\mathcal{B} = \{(x_i,\,y_i)\}_{i=1}^{N}$};

\node[attackblk, below=of batch] (attacks)
  {\textbf{Mixed Source-Attack Batch}\\[3pt]
   split minibatch across
   $\mathrm{PGD}\text{-}\ell_\infty$ and $\mathrm{DDN}\text{-}\ell_2$
   \quad$\longrightarrow$\quad
   $\{x_i^{\mathrm{adv}}\}_{i=1}^{N}$};

\node[fwdblk, below=of attacks] (forward)
  {\textbf{Model Signals from } $f_\theta(x_i^{\mathrm{adv}})$\\[3pt]
   forward: logits $z_i$, representation $h_i$
   \quad$|$\quad
   gradient: fingerprint $g_i$};

\node[featblk, below=of forward] (augfeat)
  {\textbf{Augmented Clustering Feature}\\[3pt]
   $\psi_i = \bigl[\,h_i \;;\; z_i \;;\; y_i \;;\; g_i\,\bigr]$};

\node[clustblk, below=of augfeat] (cluster)
  {\textbf{Cluster Assignment}\enspace (current $k$-means)\\[3pt]
   $c_i \;\in\; \{1,\dots,K\}$};

\node[lossblk] (lerm) at ($(cluster.south)+(-2.8,-1.95)$)
  {$\mathcal{L}_{\mathrm{ERM}}$\\[6pt]
   {\footnotesize uniform multi-attack}};

\node[lossblk] (ldro) at ($(cluster.south)+(+2.8,-1.95)$)
  {$\mathcal{L}_{\mathrm{ClusterDRO}}$\\[6pt]
   $\displaystyle\sum_{k=1}^{K} q_k\,\mathcal{L}_k$};

\coordinate (lossmid) at ($(lerm.south)!0.5!(ldro.south)$);

\node[combblk] (combine) at ($(lossmid)+(0,-1.35)$)
  {$\mathcal{L}_{\mathrm{total}}
    \;=\;
    (1-\lambda_{\mathrm{DRO}})\,\mathcal{L}_{\mathrm{ERM}}
    \;+\;
    \lambda_{\mathrm{DRO}}\,\mathcal{L}_{\mathrm{ClusterDRO}}$};

\node[updblk, below=0.65cm of combine] (updatetheta)
  {\textbf{Update Model Parameters}\\[3pt]
   $\theta \;\leftarrow\; \theta
      - \eta_\theta\,\nabla_\theta\,\mathcal{L}_{\mathrm{total}}$};

\node[updblk, below=0.65cm of updatetheta] (updateq)
  {\textbf{Update Cluster Weights}\enspace
   (EG update $+$ normalisation $+$ $q$\text{-floor})\\[3pt]
   $q_k \;\leftarrow\;
     \mathrm{Proj}_{q_{\min}}\!\!\left(
       q_k \cdot \exp\!\bigl(\eta_q\,\mathcal{L}_k\bigr)\right)$};

\node[sideblk] (bank) at ($(augfeat.east)+(3.75,0)$)
  {\textbf{Feature Bank}\\[3pt]
   $\{\psi_i\}$ \;(detached)};

\node[sideblk, below=1.50cm of bank] (kmeans)
  {\textbf{$k$-Means Clusterer}\\[3pt]
   $K$ centroids};

\draw[arr] (batch)     -- (attacks);
\draw[arr] (attacks)   -- (forward);
\draw[arr] (forward)   -- (augfeat);
\draw[arr] (augfeat)   -- (cluster);

\coordinate (splitpt) at ($(cluster.south)+(0,-0.44)$);
\draw[arr] (cluster.south) -- (splitpt);
\draw[arr] (splitpt) -| (lerm.north);
\draw[arr] (splitpt) -| (ldro.north);

% Merge the two loss branches before entering the total-loss box.
\coordinate (lossJoinR) at ($(ldro.south)+(0,-0.32)$);
\coordinate (lossJoinL) at (lerm.south |- lossJoinR);
\coordinate (lossJoinM) at ($(lossJoinL)!0.5!(lossJoinR)$);
\draw[line width=0.95pt, color=black!55] (lerm.south) -- (lossJoinL);
\draw[line width=0.95pt, color=black!55] (ldro.south) -- (lossJoinR);
\draw[line width=0.95pt, color=black!55] (lossJoinL) -- (lossJoinR);
\draw[arr] (lossJoinM) -- (combine.north);

\draw[arr] (combine)     -- (updatetheta);
\draw[arr] (updatetheta) -- (updateq);

\draw[darr] (augfeat.east)
    to node[above, font=\scriptsize\itshape, color=gray!65]
       {store (detach)}
    (bank.west);

\draw[darr] (bank.south)
    to node[right, font=\scriptsize\itshape, color=gray!65,
            xshift=2pt, align=left]
       {periodic\\[-2pt]refresh\\[-2pt]($R$ epochs)}
    (kmeans.north);

\draw[darr] (kmeans.west) -- ++(-0.65,0)
    |- node[left, font=\scriptsize\itshape, color=gray!65,
            xshift=-2pt, near start] {assign}
    (cluster.east);

\coordinate (loopbot)  at ($(updateq.south)+(0,-0.58)$);
\coordinate (loopleft) at ($(loopbot)+(-6.80,0)$);
\draw[looparr]
    (updateq.south) -- (loopbot) -- (loopleft) |- (batch.west);

\node[font=\scriptsize\itshape, color=black!32,
      anchor=north, inner sep=2pt]
    at ($(loopbot)!0.5!(loopleft)$)
    {next iteration};

\end{tikzpicture}%
}
\caption[Core AttackDRO++ training loop]{%
  \textbf{Core training loop of AttackDRO++.}
  At each iteration, the minibatch is split across two source attacks
  ($\mathrm{PGD}\text{-}\ell_\infty$ and $\mathrm{DDN}\text{-}\ell_2$),
  producing a mixed adversarial batch.
  The batch is passed through \(f_\theta\) to obtain logits \(z_i\),
  penultimate representations \(h_i\), and gradient fingerprints \(g_i\).
  These signals are concatenated into augmented features
  \(\psi_i=[h_i;z_i;y_i;g_i]\) and assigned to one of \(K\) latent clusters by
  the current \(k\)-means clusterer.
  The objective interpolates a uniform multi-attack ERM term
  \(\mathcal{L}_{\mathrm{ERM}}\) with a cluster-DRO term
  \(\mathcal{L}_{\mathrm{ClusterDRO}} = \sum_k q_k\mathcal{L}_k\) using
  anchor coefficient \(\lambda_{\mathrm{DRO}}\).
  Model parameters \(\theta\) are updated every step, while cluster weights
  \(\mathbf{q}\) are updated after the warmup period using exponentiated-gradient
  updates with floor projection; the detached feature bank periodically refreshes
  the \(k\)-means centroids every \(R\) epochs.%
}
\label{fig:attackdropp_core_training}
\end{figure}


\begin{algorithm}[H]
\caption{AttackDRO++ with Gradient Fingerprints and Uniform Anchor}
\label{alg:ch4_attackdropp}
\DontPrintSemicolon
\SetAlgoLined

\KwIn{Training set \(\mathcal{D}_{\mathrm{train}}\), model \(f_\theta\), source attacks \(\mathcal{A}_{\mathrm{src}}=\{A_1,\ldots,A_M\}\), Cluster count \(K\), anchor weight \(\lambda_{\mathrm{DRO}}\), learning rates \(\eta_\theta,\eta_q\), \(q\)-floor \(q_{\min}\), Label weight \(\lambda_{\mathrm{lbl}}\), gradient weight \(\lambda_{\mathrm{grad}}\), projection dimension \(d_{\mathrm{proj}}\), Refresh interval \(R\), warmup epochs \(T_w\)}
\KwOut{Trained model \(f_\theta^\ast\)}

Initialize \(\mathbf{q} \leftarrow (1/K,\ldots,1/K)\)\;
Initialize feature bank \(\mathcal{F} \leftarrow \emptyset\)\;
Initialize K-means clusterer as unavailable\;
Sample fixed projection matrix \(P \in \mathbb{R}^{C d_h \times d_{\mathrm{proj}}}\)\;

\For{epoch \(e = 1,\ldots,E\)}{
    \For{each minibatch \((x,y)\) of size \(N\)}{
        Generate mixed adversarial batch \(x^{\mathrm{adv}}\) using \(\mathcal{A}_{\mathrm{src}}\)\;
        Compute logits and features \((z,h) \leftarrow f_\theta(x^{\mathrm{adv}})\)\;
        Build augmented features \(\psi \leftarrow \mathrm{BuildFeature}(z,h,y,P)\)\;
        Assign cluster ids \(c \leftarrow \mathrm{AssignCluster}(\psi,\mathrm{Clusterer},K)\)\;

        Compute uniform loss \(\mathcal{L}_{\mathrm{ERM}}\)\;
        Compute cluster-DRO loss \(\mathcal{L}_{\mathrm{ClusterDRO}}\)\;
        Compute final loss
        \[
        \mathcal{L}
        =
        (1-\lambda_{\mathrm{DRO}})\mathcal{L}_{\mathrm{ERM}}
        +
        \lambda_{\mathrm{DRO}}\mathcal{L}_{\mathrm{ClusterDRO}} .
        \]

        Update model parameters \(\theta\) using SGD\;

        \If{\(e > T_w\)}{
            Update \(\mathbf{q}\) using exponentiated gradient and floor projection\;
        }

        Add detached features \(\psi\) to feature bank \(\mathcal{F}\)\;
    }

    \If{refresh condition is met}{
        Refresh K-means clusterer using \(\mathrm{KMeans}(\mathcal{F},K)\)\;
        Apply \(q\)-refresh policy to \(\mathbf{q}\) \tcp*{default: reset to uniform}
        Reset or update feature bank \(\mathcal{F}\) according to the memory policy\;
    }
}

\Return{\(f_\theta^\ast\)}
\end{algorithm}


\begin{table}[H]
\centering
\caption{Default hyperparameters for AttackDRO++ in the main experiments.}
\label{tab:ch4_hparams}
\small
\begin{tabularx}{\textwidth}{@{}l l l X@{}}
\toprule
Hyperparameter & Symbol & Value & Role \\
\midrule
DRO anchor strength      & \(\lambda_{\mathrm{DRO}}\)  & 0.35  & Interpolates between uniform ERM and Cluster-DRO \\
Number of clusters       & \(K\)                       & 4     & Controls latent group granularity \\
\(q\) learning rate      & \(\eta_q\)                  & \(3\times10^{-4}\) & Step size for exponentiated-gradient update \\
\(q\)-floor              & \(q_{\min}\)                & 0.03  & Minimum mass retained by each cluster \\
\(q\)-warmup             & \(T_w\)                     & 3 epochs & Delays adaptive reweighting until clusters are more reliable \\
Gradient projection dim. & \(d_{\mathrm{proj}}\)       & 128   & Dimension of projected gradient fingerprint \\
Gradient weight          & \(\lambda_{\mathrm{grad}}\) & 0.1   & Weight of GradFP component in \(\psi\) \\
Label weight             & \(\lambda_{\mathrm{lbl}}\)  & 0.1   & Weight of label component in \(\psi\) \\
Feature bank size        & \(|\mathcal{F}|\)           & 4{,}096 & Buffer used to refit K-means \\
Cluster refresh interval & \(R\)                       & 2 epochs & Default K-means refit frequency \\
\(q\)-refresh policy     & N/A                         & reset & Resets \(\mathbf{q}\) to uniform after cluster refresh \\
Bootstrap assignment     & N/A                         & random & Used before the first K-means refit \\
\bottomrule
\end{tabularx}
\end{table}

This chapter formulated cross-attack adversarial training as a domain-aware
optimization problem and introduced the progression from Multi-AT to AttackDRO
and AttackDRO++. The final method combines cluster-based grouping, gradient
fingerprints, and a uniform anchor to emphasize difficult adversarial groups
while retaining the stability of the uniform multi-attack baseline.

The chapter has defined the method and its training pipeline, but it has not yet
established how the method behaves empirically. Chapter~\ref{chap:experimental_setup}
therefore specifies the dataset, attacks, metrics, and evaluation protocols used
to evaluate the proposed method.

